\section{Two raw-kernel realizations}\label{sec:realizations}
The next results verify the preceding general theorem from physical kernels. Their finite affine representations are derived here and are not used as premises in Sections~\ref{sec:occupation-proof}--\ref{sec:profile-proof}.

\subsection{Correlated controlled geometric observations}
Fix $d\ge1$, let $K\subset[-1/2,1/2]^d$ be compact, and let $\nu$ be a centered probability with support $K$. The physical state space is the whole box; preparation at the origin is allowed even when the origin is outside $K$. A finite set of controls has unknown, fixed coefficients $\theta_a\in[-\eta,\eta]$, where $0\le\eta<1$. One raw acquisition from physical state $x$ selects the new state $y$ and reports it exactly, with kernel
\begin{equation}\label{eq:classical-kernel}
 K_a^\theta(dx',dy\mid x)
 =k_a^\theta(x,y)\,\delta_y(dx')\nu(dy),\qquad
 k_a^\theta(x,y)=1+\frac{4\theta_a}{d}\ip{x}{y}.
\end{equation}
Centering normalizes the kernel and $1-\eta\le k_a^\theta\le1+\eta$. The next report is generally correlated with the previous physical state. Different control coefficients change its conditional law. No reset to $\nu$ is made after a rejection. Preparation at $x=0$ produces $\nu$ in one actual call, uniformly in $\theta$.

The protected audit has conditionally independent coordinates $Z_i\in\{-1,0,1\}$, with
\begin{equation}\label{eq:classical-audit}
 \Pp(Z_i=\pm1\mid x)=\frac{x_i^2+1/2\pm x_i}{2},
 \qquad \Pp(Z_i=0\mid x)=\frac12-x_i^2.
\end{equation}
All these probabilities are positive on the box. Direct calculation gives $\E[Z_i\mid x]=x_i$ and $\operatorname{Var}(Z_i\mid x)=1/2$.

\begin{theorem}[Classical realization]\label{thm:classical}
The prepared kernels \eqref{eq:classical-kernel}--\eqref{eq:classical-audit} satisfy the finite affine and model-independent hypotheses of Theorem~\ref{thm:main}(i), and Definition~\ref{def:recoverable} with $a_*=1-\eta$, $b_*=1+\eta$ and $B=d/2$. If $\nu$ satisfies \eqref{eq:regularity}, then the same theorem yields, for the entire controlled unknown-coefficient family,
\begin{equation}\label{eq:classical-rate}
 \calR_\aut^\Theta(N,M)-\frac d2
       \asymp \bigl(N(M-1)\bigr)^{-2/s}.
\end{equation}
For uniform measure on the box, $s=d$. Every mean-budget lower holds for arbitrary history-dependent sensing and randomized stopping, not only for the implementing fixed action.
\end{theorem}
\begin{proof}
Let $\phi(x)=(x^\alpha)_{\alpha\in\{0,1,2\}^d}$, with the constant coordinate first, and let $D=\conv\phi([-1/2,1/2]^d)$. Work on its linear span with cone generated by $D$ and mass equal to the constant coordinate. The base is compact; the cone is closed and pointed. With $z_{e_i}$ denoting a first-moment coordinate, define the raw instrument
\begin{equation}\label{eq:classical-instrument}
 L_a^\theta(y)z=
   \left(z_0+\frac{4\theta_a}{d}\sum_i y_i z_{e_i}\right)\phi(y).
\end{equation}
It is positive on the cone because its scalar coefficient is the integral of the positive density \eqref{eq:classical-kernel} under a positive mixture of physical states. It is integrable and normalized. Its normalized output is $\phi(y)$, independently of the incoming predictive mixture. This verifies both finite affine update closure and report recoverability.

Every audit event has a probability polynomial of degree at most two in each coordinate, hence an affine evaluation on $D$. Conversely, the three univariate audit probabilities span $1,x_i,x_i^2$. Their products show that the finite audit-event family separates every coordinate of $D$. Future finite report-and-action words are obtained by positive compositions and integrations of \eqref{eq:classical-instrument}. Thus this is a finite closure of executable tests, not an arbitrary feature matrix. The task map extracts the first moments, and its second-moment evaluation is
\[
 c(\phi(x))=\frac d2+\sum_i x_i^2.
\]
On the actually acquired output $\phi(y)$, the conditional Bayes risk is $d/2$.

The raw likelihood bound $1-\eta\le k\le1+\eta$ survives conditioning on any history, since conditional mixtures preserve these measure inequalities. It holds for every adaptive choice of the next action. The first-call reference law is exactly $\nu$, and every fixed action is suitable for the upper construction. For completeness, the report kernel is continuous in the physical input, with
\[
 \norm{K_a^\theta(x,\cdot)-K_a^\theta(x',\cdot)}_\TV
 \le\frac{2\eta}{d}\int|\ip{x-x'}y|\,\nu(dy).
\]
On the full convex base the analogous bound uses the difference of first moments. Same-report predictive updates have zero discrepancy after normalization. These continuity and stability statements are consequences of the raw formula; neither a Fisher metric nor an assumed covariance geometry enters. Theorem~\ref{thm:main} now gives \eqref{eq:classical-rate}.
\end{proof}

\subsection{Noncommuting covariant quantum instruments}
Let $q\ge2$. A physical state is a density operator on $\mathbb C^q$. Let $\nu$ be the unitarily invariant probability on rank-one projections $P_\psi$. It is a probability on complex projective space, not on vectors with an observable phase. Each control selects a channel
\begin{equation}\label{eq:quantum-channel}
 \mathcal T_a^\theta(\rho)=\eta_a^\theta\frac{I}{q}\tr\rho
            +(1-\eta_a^\theta)U_a^\theta\rho(U_a^\theta)^*,
 \qquad \eta_a^\theta\in[\eta_0,1],\quad \eta_0>0,
\end{equation}
where the fixed coefficients and unitaries may be unknown. The raw instrument is
\begin{equation}\label{eq:quantum-instrument}
 L_a^\theta(d\psi)\rho
   =qP_\psi\mathcal T_a^\theta(\rho)P_\psi\,\nu(d\psi)
   =q\tr(P_\psi\mathcal T_a^\theta(\rho))P_\psi\,\nu(d\psi).
\end{equation}
It is a completely positive instrument: composition with $\mathcal T_a^\theta$ is followed by the positive measurement map $\rho\mapsto qP_\psi\rho P_\psi$. Since $\int P_\psi\,d\nu=I/q$, its integrated trace is preserved. The report is $P_\psi$ and the normalized output state is $P_\psi$.

After stopping, a protected covariant measurement on that output state reports $P_\zeta$ with density $q\tr(\rho P_\zeta)$ relative to $\nu$. The numerical audit vector, regarded in the real Hilbert space of Hermitian matrices, is
\begin{equation}\label{eq:quantum-audit}
 Z=(q+1)P_\zeta-I.
\end{equation}
The task is its Hilbert--Schmidt squared prediction loss. No further copy or quantum memory is available to the observer.

\begin{theorem}[Quantum realization]\label{thm:quantum}
The prepared instrument \eqref{eq:quantum-instrument} and audit \eqref{eq:quantum-audit} verify Theorem~\ref{thm:main}(i)--(iii), with
\[
 a_*=\eta_0,\qquad b_*=q,\qquad B=q^2+q-2,
 \qquad s=2(q-1).
\]
Preparation at $I/q$ realizes the reference law $\nu$ in one call for every model. Uniformly over the allowed channel family,
\begin{equation}\label{eq:quantum-rate}
 \calR_\aut^\Theta(N,M)-(q^2+q-2)
       \asymp \bigl(N(M-1)\bigr)^{-1/(q-1)}.
\end{equation}
The instrument has continuously many, generally noncommuting, state-dependent effects. No shared eigenbasis or finite designed measurement ensemble is required.
\end{theorem}
\begin{proof}
Use the real space of Hermitian matrices, its positive semidefinite cone and trace as mass. Positivity, compactness of the density-operator base and finite-dimensional affine closure follow directly from \eqref{eq:quantum-instrument}. The likelihood obeys
\[
 \eta_0\le q\tr(P_\psi\mathcal T_a^\theta(\rho))\le q
\]
for every density operator, hence for every incoming posterior mixture and adaptive history. The normalized map is the reported $P_\psi$, so it is report-recoverable. The action family changes probabilities, not the observer's access to the model.

We give the audit and geometry calculations. Unitary invariance, or the fourth moments of a normalized complex Gaussian vector, gives
\begin{equation}\label{eq:haar-second}
 \int P_\zeta\otimes P_\zeta\,\nu(d\zeta)
       =\frac{I+\mathsf{Swap}}{q(q+1)}.
\end{equation}
To verify this formula directly, independent standard complex Gaussian coordinates have independent exponential squared moduli. Their normalized squared moduli have uniform simplex density, with second moments $2/(q(q+1))$ on the diagonal and $1/(q(q+1))$ off the diagonal; independent phases annihilate the remaining entries. Pairing \eqref{eq:haar-second} with $\rho$ yields
\[
 \E[P_\zeta\mid\rho]=\frac{\rho+I}{q+1},\qquad \E[Z\mid\rho]=\rho.
\]
Moreover, $Z$ has eigenvalues $q,-1,\ldots,-1$, so $\norm Z_\HS^2=q^2+q-1$. On each acquired pure state, $\norm\rho_\HS^2=1$, giving constant Bayes risk $B=q^2+q-2$. The audit separates all density operators, since its mean is $\rho$. It therefore supplies executable informationally complete tests, rather than an assumed state metric.

For a fixed ray $P_\psi$, put $V=|\ip\psi\zeta|^2$. The same exponential-coordinate calculation gives
\[
 \Pp(V\ge u)=(1-u)^{q-1},\quad 0\le u\le1.
\]
Since $\norm{P_\psi-P_\zeta}_\HS^2=2(1-V)$, the actual reference measure has the exact ball mass
\begin{equation}\label{eq:projective-mass}
 \nu\{P_\zeta:\norm{P_\psi-P_\zeta}_\HS\le r\}
       =\left(\frac{r^2}{2}\right)^{q-1},\quad 0\le r\le\sqrt2.
\end{equation}
This proves the required regularity on the acquired projective support. In particular, its dimension is $2(q-1)$, not the ambient real dimension $q^2$ of Hermitian coordinates.

Finally, trace-norm contraction of a channel and the elementary bound $|\tr(PA)|\le\norm A_1$ give report-law continuity,
\[
 \norm{p_{a,\rho}^\theta-p_{a,\sigma}^\theta}_\TV
 \le\frac q2\norm{\rho-\sigma}_1.
\]
The posterior for a common report is exactly the same projection. Thus the raw likelihood, actual mass, global cover, update stability and protected audit conditions have all been verified. Apply Theorem~\ref{thm:main}.
\end{proof}

\begin{remark}[The role of physical preparation]
The quantum output projection and the classical next state are states of the experiment, not extra registers of the observer. The observer retains only its current classical label. A call to the quantum instrument includes the displayed controlled channel and measurement; the number of such calls is $\tau$. Gate counts, channel preparation, the protected audit call and the fixed apparatus dimension are additional physical costs. They are not equated with $\lceil\log_2M\rceil$ classical bits. The rate is a statement for this declared instrument interface, not for arbitrary laboratory implementations with free resets.
\end{remark}
